% academic-paper — arxiv generic scaffold
%
% Open-license starter for arxiv preprints. For NeurIPS/ICML/ICLR/CHI,
% download the venue's official template instead (see POINTERS.md).
%
% Usage:
%   1. Copy this file to your paper directory. Rename (e.g., paper.tex).
%   2. Create references.bib for your citations.
%   3. Compile: pdflatex paper.tex, bibtex paper, pdflatex paper.tex, pdflatex paper.tex
%   4. For arxiv upload: ensure paper.bbl exists and shares the base name with paper.tex.

\documentclass[11pt]{article}

\usepackage[utf8]{inputenc}
\usepackage[T1]{fontenc}
\usepackage[letterpaper,margin=1in]{geometry}
\usepackage{natbib}
\usepackage{graphicx}
\usepackage{booktabs}
\usepackage{hyperref}
\usepackage{amsmath}
\usepackage{xcolor}
\usepackage{enumitem}

\hypersetup{
  colorlinks=true,
  linkcolor=black,
  citecolor=blue,
  urlcolor=blue,
}

\bibliographystyle{plainnat}

% Paper metadata
\title{Your Paper Title: Subtitle If Needed}

\author{%
  First A. Author\thanks{We use authorial ``we'' throughout, following CS convention.} \\
  Affiliation \\
  City, Country \\
  \texttt{email@example.com}
  \and
  Second B. Author \\
  Affiliation \\
  \texttt{email@example.com}
}

\date{}

\begin{document}

\maketitle

\begin{abstract}
% Four-move formula:
% 1. Problem (1-2 sentences)
% 2. Approach (1-2 sentences)
% 3. Key results with specific numbers (2-3 sentences)
% 4. Significance (1 sentence)
% Target: 150-250 words.
\end{abstract}

\section{Introduction}
% CARS moves:
% Move 1: Establish territory (why the field or problem matters)
% Move 2: Establish niche (the gap)
% Move 3: Occupy niche (what this paper does)
% End with contribution bullets:
%   1. We present...
%   2. We demonstrate...
%   3. We report...
%   4. We propose...

\section{Related Work}
% Organize by theme, not by paper.
% Each subsection: what the category does, what it does not do, how this paper differs.

\subsection{Category 1}

\subsection{Category 2}

\section{Method}
% For empirical: IMRaD Methods.
% For theory: Formal definitions, argument structure.

\subsection{Participants or Dataset}

\subsection{Procedure}

\subsection{Analysis}

\section{Results}
% Present findings. Minimal interpretation (save for Discussion).
% Active voice. Past tense for what was observed.
% Include effect sizes, CIs, and p-values for quantitative.
% Include thick description and quotes for qualitative.

\section{Discussion}
% Summarize findings. Interpret. Connect to prior work.
% Heaviest citation density. Appropriate hedging.

\section{Limitations}
% Honest accounting. One limitation per sentence or paragraph.
% No softeners. No boilerplate.

\section{Conclusion}
% Restate thesis with evidence. Short.
% No new information.

% Example figure
% \begin{figure}[htbp]
%   \centering
%   \includegraphics[width=0.8\textwidth]{figures/diagram.pdf}
%   \caption{Caption that stands alone with the figure.}
%   \label{fig:diagram}
% \end{figure}

% Example table (with booktabs)
% \begin{table}[htbp]
%   \centering
%   \begin{tabular}{lrr}
%     \toprule
%     Metric & Before & After \\
%     \midrule
%     Value 1 & 0.42 & 0.87 \\
%     Value 2 & 3.14 & 2.71 \\
%     \bottomrule
%   \end{tabular}
%   \caption{Table caption.}
%   \label{tab:results}
% \end{table}

\bibliography{references}

% Optional appendix
% \appendix
% \section{Supplementary Material}

\end{document}
